\subsection{Comparação das versões}

A Tabela~\ref{tab:final} mostra as versões nos 5.842 looks de teste. Os arquivos de origem de cada número estão no repositório~\cite{repositorio}. Todos os valores de Hit@k são baixos porque cada consulta tem, em média, 10.940 candidatas. Por isso o sorteio acerta só 0,182\% no Hit@10. A versão final F chega a 1,056\% no Hit@10 e a 32,04\% no FITB. A popularidade tem 0,895\% e 31,40\%, e a versão C tem 0,751\% e 25,37\%.

A Tabela~\ref{tab:diferencas} separa o que é diferença real do que pode ser sorte. Contra a C, a F ganha no Hit@10 por 0,305 ponto percentual (intervalo de 0,113 a 0,510) e no FITB por 6,67 pontos (5,90 a 7,43). Os dois intervalos ficam acima de zero, então a vantagem é real. Contra a popularidade, a F ganha no FITB por 0,63 ponto (0,14 a 1,11). No Hit@10, porém, o intervalo vai de $-$0,068 a 0,383 e a diferença pode ser sorte. Na validação, a F havia ganhado da popularidade no Hit@10. Mas a F foi escolhida nesses mesmos looks, e isso favorece a versão escolhida. O teste não confirmou essa vantagem. Contra a C + popularidade, o Hit@1 é igual e os intervalos das outras métricas passam pelo zero. Portanto, o ganho da F sobre a C vem da popularidade. Os grupos do Kruskal não trazem ganho de acerto que se possa medir. Na validação, eles venceram por uma única consulta em 2.332. Ficaram na F pelo critério definido antes da rodada.

A F perde no Hit@1 para a B e a C (0,313\% contra 0,373\% e 0,336\%) e no FITB para a A (34,29\%). A A mede só a semelhança de atributos com o look. Ela separa bem a peça certa de três sorteadas, mas raramente a coloca entre as 10 primeiras (0,514\%). A B e a C ficam no nível do sorteio no FITB (25,35\% e 25,37\%). O motivo é que, na maior parte das consultas, a escondida tem nota zero. Ela então empata, em média, com 10.880 peças na B e 10.870 na C.

\subsection{Peça conhecida e peça nova}

A Tabela~\ref{tab:conhecida_nova} divide o teste em dois grupos. Em 1.600 consultas (27,4\%), a peça escondida aparece em algum look de treino. Nas outras 4.242 (72,6\%), ela é nova. Na peça conhecida, a F acerta 3,856\% no Hit@10, 23,0 vezes o aleatório. Na peça nova, o Hit@10 da F é 0,000\%, e todo o acerto dela vem da peça conhecida. A causa está na estrutura do grafo. A peça nova tem grau zero no grafo de peças e looks, então a BFS nunca chega até ela. A popularidade, $\log(1+\text{grau})$, também vale zero. Assim, qualquer peça conhecida da categoria fica à frente. Nesse grupo, só as versões que não dependem dos looks superam o aleatório. A A chega a 0,542\% (2,9 vezes o aleatório) e a 35,10\% no FITB, com vantagem real sobre a F nas duas métricas. Os atributos do texto trazem alguma informação sobre a peça nova. Mas eles não competem com o sinal dos looks na peça conhecida.

\subsection{Tempos: teoria e prática}

A Tabela~\ref{tab:tempos} e a Fig.~\ref{fig:tempos} trazem a média de 5 repetições em subconjuntos reais da base. A construção da KD-tree cresce perto de $n \log n$. Em $n = 10^5$, o kNN na KD-tree leva 0,09484 ms por consulta, contra 1,343 ms da força bruta, ou seja, 14 vezes menos. Em $n = 10^3$, a diferença é de só 1,6 vez. Porém a busca na KD-tree ficou 13 vezes mais lenta de $10^3$ para $10^5$, enquanto $\log n$ cresce menos de 2 vezes. A razão é que muitas peças têm os mesmos valores de formalidade e de cor. Com pontos repetidos, a busca visita mais ramos do que prevê o caso médio~\cite{friedman1977}.

No Union-Find sem união por tamanho e sem compressão de caminho, a sequência de teste forma uma corrente de $n$ elementos. Quando $n$ fica 10 vezes maior, o tempo fica perto de 100 vezes maior (3,458 ms, 297,6 ms e 29.204 ms), como prevê o custo quadrático. Com as duas técnicas~\cite{galler1964,tarjan1975}, o crescimento é linear. Em $n = 10^5$, o tempo cai para 0,7813 ms, cerca de 37 mil vezes menos. No Kruskal, ordenar as arestas custa de 7,1 a 8,7 vezes o laço de uniões, como prevê o custo $O(E \log E)$. O Kruskal foi de 1,4 a 1,8 vez mais rápido que o Prim com heap. As duas florestas tiveram o mesmo peso. A BFS de dois passos leva menos de 1 microssegundo por consulta. Ela depende do tamanho dos looks da peça, e não de $n$.

O gargalo é a montagem do grafo kNN usado pelo Kruskal. Ela faz uma busca kNN por peça e leva 8.012 ms em $n = 10^5$. Ela cresce mais que $n \log n$ pelo mesmo motivo da busca na KD-tree. Com a base maior, as etapas feitas uma vez crescem perto de $n \log n$, exceto pelo efeito dos valores repetidos. Já cada recomendação cresce de forma linear com o tamanho da categoria, pois dá nota a todas as peças dela. Uma rodada com a máquina mais carregada deu tempos de 10\% a 40\% maiores, sem mudar a forma das curvas.

\subsection{Grupos e articulações}

A Fig.~\ref{fig:grupos} mostra o grafo entre grupos, desenhado no Gephi~\cite{bastian2009} com o layout ForceAtlas2~\cite{jacomy2014}. Os vértices são os 1.255 grupos de cor, estampa e formalidade dentro de cada categoria. Não usamos os grupos do Kruskal, pois neles o efeito corrente deixa 16 de 20 categorias com um grupo de mais da metade das peças. Uma aresta liga grupos de categorias diferentes com pelo menos 3 pares de peças em comum e pontos acima da mediana (0,437). O tamanho do vértice indica o número de peças. A cor destaca as articulações que separam 3 grupos ou mais. Há 711 grupos ligados por 5.604 arestas em um único componente.

O algoritmo de Tarjan~\cite{tarjan1972,hopcroft1973} encontrou 53 articulações e 96 pontes em 126~$\mu$s. Só 6 articulações separam 3 grupos ou mais. A principal é o grupo de tops brancos lisos de formalidade baixa, que separa 10 grupos. Em seguida vêm os grupos de sapatos pretos lisos e de bolsas pretas lisas de formalidade alta, com 8 cada. Peças básicas e neutras são a única ligação de alguns grupos com o resto do grafo. Porém, em todas as articulações, a maior parte que se separa tem um único grupo. Nenhuma ponte separa mais de 1 grupo. Ou seja, as articulações prendem grupos isolados na borda do grafo, e não estilos inteiros. Grau e articulação também medem coisas diferentes. O grupo de maior grau, com 96 vizinhos, separa só 2 grupos.

\subsection{Experimento registrado: a peça nova}

Depois de fixar a F, testamos 596 configurações para a peça nova. A escolha usou só a validação, dividida em duas metades: a metade A escolhe e a metade B confirma. Pelo Hit@10, nenhuma configuração ganhou da F nas duas metades. A melhor na metade A chegou a 1,630\%, contra 1,123\% da F. Na metade B, ela caiu para 0,995\%, contra 1,312\% da F. Pelo FITB, critério acrescentado depois, a melhor variação foi ao teste uma única vez, com o nome F2. Ela chegou a 41,64\% no FITB, uma diferença real de 9,60 pontos sobre a F (8,48 a 10,70), e a 37,15\% na peça nova. No Hit@10, a diferença para a F pode ser sorte, e na peça nova o valor continua em 0,000\%. A F2 fica fora das tabelas, pois o critério principal era o Hit@10. Os atributos e os grupos ajudam a escolher entre poucas opções. Mas não colocam a peça nova à frente das conhecidas na categoria inteira.
